\documentclass{article}
\usepackage{amsmath}
\usepackage{amssymb}
\usepackage{graphicx}
\usepackage{accessibility}

\title{A Computational Framework for Exploring Algorithmic Information Bounds in Geometric Measure Theory}
\author{Jules}
\date{\today}

\begin{document}

\maketitle

\begin{abstract}
This paper introduces an open-source Rust library for the computational exploration of algorithmic and probabilistic information theory alongside universal register machines, drawing on "Algorithmic Information Bounds for Distances and Orthogonal Projections" by Cholak et al.\ together with Conway's FRACTRAN universal arithmetic language and classical Shannon-Blahut-Arimoto channel coding theory. We present the design of the library, which includes modules for geometric primitives, Kolmogorov complexity approximations, the core Combinatorial Lemma, a zero-allocation prime-exponent FRACTRAN virtual machine with Minsky counter-machine decoding, and preallocated discrete Shannon entropy, divergence, and channel capacity solvers. We demonstrate the library's utility through numerical experiments on Hausdorff dimension bounds, prime-generating counter automata, and discrete memoryless channel capacities.
\end{abstract}

\section{Introduction}
% Overview of the problem
% Motivation for the work
% Summary of the paper's structure

\section{Mathematical Preliminaries}
% Review of Kolmogorov complexity
% Review of Hausdorff dimension
% Review of the point-to-set principle
% Statement of the key theorems from the source paper

\section{Implementation Design}

The core of our computational framework is a Rust library exposed through the \texttt{math\_explorer} crate. The library is designed to be modular and extensible, with a clear separation of concerns that reflects the underlying computability, probabilistic information, and geometric measure concepts. The main modules are \texttt{geometry}, \texttt{kolmogorov}, \texttt{combinatorics}, \texttt{fractran}, and \texttt{shannon}.

\subsection{The \texttt{geometry} Module}

The \texttt{geometry} module provides the foundational data structures for representing geometric objects. To handle the arbitrary-precision arithmetic required by the theoretical framework, we wrap exact rational numbers (\texttt{dashu::rational::RBig}) inside \texttt{Rational}.

\begin{itemize}
    \item \texttt{Point2D}: A type alias for a 2D vector of \texttt{Rational} numbers, representing a point in $\mathbb{R}^2$.
    \item \texttt{Line}: A struct representing a line in 2D space, defined by an origin point and a direction vector.
    \item \texttt{DyadicRational}: A struct representing a dyadic rational number of the form $m/2^r$.
    \item \texttt{DyadicPoint2D}: A struct representing a point with dyadic rational coordinates.
\end{itemize}

The module also provides functions for essential geometric operations, such as calculating Euclidean distances and projecting points onto lines.

\subsection{The \texttt{kolmogorov} Module}

The \texttt{kolmogorov} module provides functions for working with computable approximations of Kolmogorov complexity. As prefix Kolmogorov complexity ($K$) is not computable, we implement functions that compute upper bounds and approximations as used in the paper.

\begin{itemize}
    \item \texttt{prefix\_kolmogorov\_approx}: This function computes an upper bound for $K(n)$ for a given integer $n$, based on the inequality $K(n) \le \log_2(n) + 2\log_2(\log_2(n)) + O(1)$.
    \item \texttt{k\_floor\_r}: This function approximates the complexity of a point at a given precision $r$ by computing the complexity of its truncated binary expansion, $K(\lfloor x \rfloor_r)$. This is a key technique used in the paper to bridge the gap between the uncomputable nature of $K$ and the need for concrete calculations.
\end{itemize}

\subsection{The \texttt{combinatorics} Module}

The \texttt{combinatorics} module implements the Combinatorial Lemma (Lemma 1) from the source paper. This lemma is the algorithmic heart of the paper's main proofs.

\begin{itemize}
    \item \texttt{combinatorial\_lemma}: This generic function takes as input the sets $X$ and $V$, the collection of subsets $N_v$, and the similarity relations $\sim_d$, and it performs a search to find a pair of points $(u, v)$ that satisfies the lemma's conclusion. While the lemma's proof is a non-constructive counting argument, our implementation provides a constructive way to find such a pair, albeit with a computational cost that depends on the size of the sets.
\end{itemize}

\subsection{The \texttt{fractran} Module}

The \texttt{fractran} module implements John Horton Conway's FRACTRAN universal register language by compiling positive rational multipliers $f_j = A_j / B_j$ into preallocated prime-exponent tables over a sorted prime basis $p_0 < p_1 < \dots < p_{m-1}$. Representing the machine integer $N = \prod_{i=0}^{m-1} p_i^{e_i}$ directly by its exponent vector $\mathbf{e} \in \mathbb{N}^m$ avoids arbitrary-precision integer blow-up and enables deterministic zero-allocation execution.

\begin{itemize}
    \item \texttt{FractranProgram} and \texttt{FractranCompiledVm}: \texttt{FractranProgram} validates strictly positive \texttt{Rational} rules over smooth prime factors and compiles their numerators and denominators into flat \texttt{u32} exponent matrices. \texttt{FractranCompiledVm::step\_registers} finds the first rule satisfying $\mathbf{e} \ge \nu_{\mathbf{p}}(B_j)$ and updates $\mathbf{e} \leftarrow (\mathbf{e} - \nu_{\mathbf{p}}(B_j)) + \nu_{\mathbf{p}}(A_j)$ in place using checked \texttt{u32} arithmetic with zero dynamic heap allocations.
    \item \texttt{FractranState} and Minsky Counter Machines: Partitioning the prime basis into indicator control-state primes and unbounded data-register primes decodes any FRACTRAN configuration into an explicit Minsky counter-machine state and extracts directed guard/increment transitions.
    \item \texttt{FractranPreset}: Built-in canonical programs include Conway's 14-fraction universal prime generator \texttt{PRIMEGAME} (emitting pure powers of two $2^2, 2^3, 2^5, 2^7, 2^{11}, \dots$), \texttt{MULT}, \texttt{ADD}, and \texttt{FIBONACCI}.
\end{itemize}

\subsection{The \texttt{shannon} Module}

To bridge individual-object Kolmogorov bounds $K(x)$ with ensemble-average coding limits, the \texttt{shannon} module provides deterministic, zero-allocation \texttt{f64} primitives (in base-$2$ bits) verified under \texttt{VerifiedAllocator}.

\begin{itemize}
    \item \texttt{categorical\_entropy} and \texttt{renyi\_entropy}: Evaluates discrete Shannon entropy $H(P) = -\sum_{i} p_i \log_2 p_i$ (with $0 \log_2 0 = 0$) and order-$\alpha$ R\'enyi entropy
    \begin{equation}
        H_{\alpha}(P) = \frac{1}{1 - \alpha} \log_2 \sum_{i} p_i^{\alpha}, \qquad \alpha \in [0, \infty],
    \end{equation}
    covering Hartley support entropy ($\alpha = 0$), the continuous Shannon limit ($\alpha \to 1$), collision entropy ($\alpha = 2$), and min-entropy $H_{\infty}(P) = -\log_2 \max_i p_i$ via log-sum-exp rescaling.
    \item \texttt{kl\_divergence}, \texttt{jensen\_shannon\_divergence}, and \texttt{joint\_mutual\_information}: Computes Kullback--Leibler relative entropy $D_{\mathrm{KL}}(P \parallel Q) = \sum_i p_i \log_2(p_i / q_i)$ with explicit absolute-continuity checking, single-pass bounded Jensen--Shannon divergence $\mathrm{JSD}(P \parallel Q) = \frac{1}{2} D_{\mathrm{KL}}(P \parallel M) + \frac{1}{2} D_{\mathrm{KL}}(Q \parallel M) \in [0, 1]$ where $M = \frac{1}{2}(P + Q)$, and joint mutual information $I(X; Y) = \sum_{x, y} P(x, y) \log_2 \frac{P(x, y)}{P_X(x) P_Y(y)}$ using preallocated marginal buffers in \texttt{JointMiWorkspace}.
    \item Closed-Form and Blahut--Arimoto Channel Capacities: Provides exact capacities for the binary symmetric channel $C_{\mathrm{BSC}}(p) = 1 - h_2(p)$, binary erasure channel $C_{\mathrm{BEC}}(\epsilon) = 1 - \epsilon$, and additive white Gaussian noise channel $C_{\mathrm{AWGN}}(\mathrm{SNR}) = \frac{1}{2}\log_2(1 + \mathrm{SNR})$ (alongside bandlimited $B \log_2(1 + \mathrm{SNR})$). For general discrete memoryless channels $W(y \mid x)$, \texttt{blahut\_arimoto\_capacity} iterates within \texttt{BlahutArimotoWorkspace} using log-stabilized updates $r_x^{(t+1)} \propto r_x^{(t)} 2^{d_x^{(t)}}$ with $d_x^{(t)} = D_{\mathrm{KL}}(W(\cdot \mid x) \parallel q^{(t)})$, bracketing $C$ between $I_L^{(t)} = \sum_x r_x^{(t)} d_x^{(t)}$ and the minimax dual bound $I_U^{(t)} = \max_x d_x^{(t)}$ without heap allocation.
\end{itemize}

This modular design allows for future extensions and provides a solid foundation for the numerical experiments and interactive visualizations presented in the next section.

\section{Numerical Experiments and Results}
% Experiments for Theorem 2 (distances)
% Experiments for Theorem 3 (projections)
% Experiments for Theorem 4 (pinned distance sets)

\section{Conclusion and Future Work}
% Summary of contributions
% Directions for future research

\bibliographystyle{plain}
\bibliography{references}

\end{document}
